The weights are constant on each embedding orbit used in the
following count. Indeed a unit of the definite order $R_z$ has
positive integral reduced norm with integral inverse, so its norm
is one. The conjugation law for $f$ therefore gives weight one
for its action. Along the class-group torsor the Schwartz factor
is constant because $a$ centralizes $x_0$, and the remaining
norm factor is exactly $\chi_F([a])$. Hence the finite-sum identity
just proved applies with the single phase
$\chi_{h_*}(\operatorname{Nrd}z_0)f(z_0^{-1}x_0z_0)$.


For each fixed $z$, the stabilizer of an embedding in $R_z^\times$
is $\mathcal O_F^\times=\{\pm1\}$. Each embedding orbit thus
has $|R_z^\times|/2$ elements. The denominator $|R_z^\times|$
in (O.4) leaves the factor $1/2$ per embedding class. The
vectors $x$ and $-x$ are not identified by this stabilizer:
at $\ell_0$ their residues define the two different embeddings
of the residue field, whereas conjugation by an order unit acts
trivially on that field. They already lie in the two different
$\ell_0$ orientation tuples counted above. The
contribution of the reference orientation tuple is
\[
 \frac{c_0}{2}\,\zeta_M(D)
                         \lambda_M(\overline{LP_D}),\qquad
 \zeta_M(D)=\chi_{h_*}(\operatorname{Nrd}z_0)
                               f(z_0^{-1}x_0z_0),
\]
up to the harmless quadratic reference-class factor just described.
All nonzero factors in this displayed phase are fixed-character
values or signs. Thus it is a $2$-power root of unity; its field
does not grow with $D$, $M$, or $\ell_0$.

It remains to account for the other orientations; they are not
all in the image of the one sign component of the reduction map.
At $p^e\Vert N$, use the local Atkin--Lehner normalizer. It
reverses the selected root, so multiplies the Schwartz weight by
\[
              \alpha_p(-1)=w_{p^e}\chi_{h_*}(p)^e.
\]
It multiplies the norm-weighted parametrization value by the same
sign: its reduced norm has valuation $e$, and the actual
parametrization transforms by $w_{p^e}$ after $L$ kills the fixed
cuspidal translation. The two signs multiply to one. This
normalizer has trivial components at the other primes, so it
changes no other orientation or local weight.

At $\ell_0$, a division-algebra uniformizer exchanges the two
residue embeddings. It sends the trace-zero residue $\beta$ to
$-\beta$, multiplying its weight by $(-1/\ell_0)=-1$. Its
reduced norm has valuation one and hence contributes
$\chi_{h_*}(\ell_0)=-1$. Its action on supersingular pairs is
Frobenius. Since $\lambda_M$ has Frobenius sign $+1$ and the
parametrization is defined over $\mathbf Q$, this introduces no
further sign. Again the two factors multiply to one.
The uniformizer operation used here may be described with arithmetic
or geometric Frobenius: they have the same action on the supersingular
pairs and their parametrization values over $\mathbb F_{\ell_0^2}$.
The identity $\lambda_M(FP)=\lambda_M(P)$ from (O.2) consequently
applies to the very operation that exchanges the residue embeddings.
Its norm has valuation one, and $\chi_{h_*}$ is unramified at
$\ell_0$, so its norm-character value is exactly
$\chi_{h_*}(\ell_0)=-1$, independently of the uniformizer's unit part.
Likewise, at $p^e\Vert N$, the norm of the Atkin--Lehner normalizer
has valuation $e$ and $\chi_{h_*}$ is unramified at $p$. Its
norm-character value is exactly $\chi_{h_*}(p)^e$. These observations
fix the norm signs in both orientation changes.


It follows that every one of the $2^{\omega(N)+1}$ tuples has
the same contribution as the reference tuple. Multiplying by
their number gives precisely $c_0 2^{\omega(N)}$, as asserted.
This counts individual embeddings with their stabilizer; a
conjugate-pair formulation of the geometric correspondence must
not be used to insert or remove another factor of two.

Finally, $c_0/|R_z^\times|$ is an integer for every $z$. Thus a
change of the chosen integer lift of $P_M(z)$ by $2^M$ changes
the actual coefficient by an element of $2^M\mathcal O$.
Every preceding equality is consequently valid at the original
retained precision. No denominator depending on the class number,
the discriminant, or the stage has entered.
\paragraph{Precision in the orientation count.}
The factor $1/2$ in the orbit count does not require dividing a
congruence modulo $2^M$ by two. Indeed, $-1\in R_z^\times$ implies
$2\mid |R_z^\times|$, and the chosen integer $c_0$ is divisible by
$|R_z^\times|$. Hence $c_0/2$ is an integer. Choose lifts of all
$P_M(z)$ before summing. Each orbit of embeddings contributes
$(c_0/|R_z^\times|)(|R_z^\times|/2)P_M(z)=(c_0/2)P_M(z)$,
with its prescribed unit weight. Replacing any lift by a multiple
of $2^M$ changes this expression by a multiple of $2^M$.
The comparison with the reduced point may therefore be made at
precision $M$, separately for each orientation. Summing the
$2^{\omega(N)+1}$ orientations gives the stated multiplier
$c_0 2^{\omega(N)}$ without any loss of precision.
In particular, write $c_\theta=9+\omega(N)$, the valuation of
the fixed integer $c_0 2^{\omega(N)}$. At the fixed witness $D_w$,
(O.2) and (O.7) imply
\[
 v_2a_M(D_w)=c_\theta+c_w
       \qquad\text{whenever }M>c_\theta+c_w.
\]
The phase is a unit and the congruence error has strictly larger
valuation, so this is an equality, not merely a lower bound.
Discarding the finitely many earlier stages is sufficient. Thus the
visibility used in choosing the odd minimum retains a fixed witness
and a fixed valuation; neither the class number nor the changing
prime $\ell_0$ introduces a new loss of precision.

\end{proof}

\begin{proposition}[Visibility and the native genus index]\label{ex:app-visibility}
